%%%%%%%%%%%%%%%%%%%%%%% CHAPTER - 7 %%%%%%%%%%%%%%%%%%%%\\
\chapter{Testing}
\label{C7} %%%%%%%%%%%%%%%%%%%%%%%%%%%%
%\graphicspath{{Figures/Chapter-6figs/PDF/}{Figures/Chapter-6figs/}}
% \noindent\rule{\linewidth}{2pt}
%%%%%%%%%%%%%%%%%%%%%%%%%%%%%%%%%%%%%%%%%%%%%%%%%%%%%%%%%%%%%%%%%%%%%%%%%%%%%%%%%%
\section{Testing} \label{S7.1}

Testing is the process used to evaluate a software product's usability, efficiency, and 
quality to make sure that it complies with the criteria and specifications laid down for 
it in the frame of software development. In order to find bugs, failures, or other 
problems that can impair the software's operation, usability, or performance, testing 
often entails running the program under a variety of circumstances and input 
circumstances. Testing is a crucial step in the development of software because it 
enables the detection and correction of any problems or flaws before the product is 
made available to consumers. Unit testing, integration testing, system testing, 
acceptance testing, and regression testing are just a few of the several forms of testing 
that may be done. Each kind of testing has a particular function and is created to handle 
particular faces of the software development process.

\subsection{Types of Testing}

Various types of testing can be carried out while developing software. Here are a few 
of the most typical: 

\textbf{\textit{Unit Testing:-}} This is the process of checking that individual software program units 
or components function as planned.

\textbf{\textit{Integration Testing:-}} This entails evaluating how various parts or units of a software 
program interact with one another to make sure they are compatible and perform as 
intended. 

\textbf{\textit{System Testing:-}} This entails testing the entire system to make sure it complies with 
the guidelines and standards established for it. 

\textbf{\textit{Acceptance Testing:-}} In preparation for a software product being made available to 
users, this is the last phase of testing. To make sure the program satisfies the user's 
needs, it must be tested in a setting similar to a production one. 

\textbf{\textit{Regression Testing:-}} Retesting a software program after modifications have been made 
is necessary to make sure that no new defects or problems were created by the changes. 

\textbf{\textit{Performance Testing:-}} This entails testing the software to make sure it operates 
effectively under real-world or hypothetical pressures. 

\textbf{\textit{Security Testing:-}} Testing the software is necessary to find any flaws or vulnerabilities 
that an attacker could use against it. 

\textbf{\textit{Usability Testing:-}} This entails evaluating the software to make sure that its intended 
users would find it simple to use and understand. 

\textbf{\textit{Compatibility Testing:-}} To make sure it works in all intended environments, this 
requires testing the software program on various hardware, software, and operating 
system configurations. 

\textbf{\textit{Exploratory Testing:-}} This entails testing the software application without following a 
predetermined test strategy with the intention of identifying fresh and unexpected flaws.

\subsection{Test Summary Report}

\begin{table}[H]
	\centering
	\caption{Decision Based Testing}
	\renewcommand{\arraystretch}{0.9}
	\begin{tabular}{|l|l|c|c|c|c|}
		\hline
		\multicolumn{2}{|l|}{\textbf{MODULE UNDER TEST: AI Cricket Coach – Shot Detection}} & 
		\multicolumn{4}{c|}{\textcolor{red}{\textbf{RULES}}} \\ \hline
		
		\multirow{3}{*}{\textbf{CONDITIONS}}
		& When a batting shot is detected 
		& \cellcolor{green!60!black}T & \cellcolor{green!60!black}T & \cellcolor{red!90}F & \cellcolor{red!90}F \\ \cline{2-6}
		
		& When the accuracy of the detection is >= 90 
		& \cellcolor{green!60!black}T & \cellcolor{red!90}F & \cellcolor{red!90}F & \cellcolor{red!90}F \\ \cline{2-6}
		
		& When stance/pose is unknown 
		& \cellcolor{red!90}F & \cellcolor{red!90}F & \cellcolor{red!90}F & \cellcolor{green!60!black}T \\ \hline
		
		\multirow{5}{*}{\textbf{ACTIONS}}
		& Display Shot Name 
		& {\color{green!60!black}\checkmark} &  &  &  \\ \cline{2-6}
		
		& Display Confidence Score 
		& {\color{green!60!black}\checkmark} & {\color{green!60!black}\checkmark} &  &  \\ \cline{2-6}
		
		& Announce Shot Name (Audio) 
		& {\color{green!60!black}\checkmark} &  &  &  \\ \cline{2-6}
		
		& Display Pose Keypoints and Skeleton Lines
		& {\color{green!60!black}\checkmark} & {\color{green!60!black}\checkmark} &  &  \\ \cline{2-6}
			
		& Display ``No Shot Detected'' 
		&  &  & {\color{green!60!black}\checkmark} &  \\ \cline{2-6}
		
		& Display ``Unknown Stance'' 
		&  &  &  & {\color{green!60!black}\checkmark} \\ \cline{2-6}
		
		& Announce ``Unknown Stance'' 
		&  &  &  & {\color{green!60!black}\checkmark} \\ \hline
		
	\end{tabular}
\end{table}

\begin{table}[H]
	\caption{Test Cases for Detection Module}
	\renewcommand{\arraystretch}{1.2}
	\setlength{\tabcolsep}{1.0pt}
	\begin{tabular}{|c|p{2.0cm}|p{2.5cm}|p{3.5cm}|p{3.5cm}|c|}
		\hline
		\rowcolor[gray]{0.85}
		\textbf{Test Case Id} & \textbf{Function Under Test} & \textbf{Input} & 
		\textbf{Expected Output} & \textbf{Actual Output} & \textbf{Pass/Fail} \\
		\hline
		
		TC-01 & Detection Module & Shot:Straight Drive &
		Display Shot Name, Display Confidence Score, Announce Shot Name, Display Pose Keypoints \& Skeleton Lines &
		Display Shot Name, Display Confidence Score, Announce Shot Name, Display Pose Keypoints \& Skeleton Lines &
		\textcolor{green!60!black}{Pass} \\
		\hline
		
		TC-02 & Detection Module & Shot:Pull Shot &
		Display Shot Name, Display Confidence Score, Announce Shot Name, Display Pose Keypoints \& Skeleton Lines &
		Display Shot Name, Display Confidence Score, Announce Shot Name, Display Pose Keypoints \& Skeleton Lines &
		\textcolor{green!60!black}{Pass} \\
		\hline
		
		TC-03 & Detection Module & Shot:Cut Shot &
		Display Shot Name, Display Confidence Score, Announce Shot Name, Display Pose Keypoints \& Skeleton Lines &
		Display Shot Name, Display Confidence Score, Announce Shot Name, Display Pose Keypoints \& Skeleton Lines &
		\textcolor{green!60!black}{Pass} \\
		\hline
		
		TC-04 & Detection Module & Shot:Sweep &
		Display Shot Name, Display Confidence Score, Announce Shot Name, Display Pose Keypoints \& Skeleton Lines &
		Display Shot Name, Display Confidence Score, Announce Shot Name, Display Pose Keypoints \& Skeleton Lines &
		\textcolor{green!60!black}{Pass} \\
		\hline
		
		TC-05 & Detection Module & Unknown Stance &
		Display ``Unknown Stance'' &
		Display ``Unknown Stance'' &
		\textcolor{green!60!black}{Pass} \\
		\hline
		
		TC-06 & Detection Module & No Person in Frame &
		Display ``No Shot Detected'' &
		Display ``No Shot Detected'' &
		\textcolor{green!60!black}{Pass} \\
		\hline
		
	\end{tabular}
\end{table}
